%----------------------------------------------------------------------------
\chapter{Further reading}\label{app:reading}
%----------------------------------------------------------------------------
\section{First-place papers to learn from}\label{app:tdk:papers}
\Cref{tab:papers-ftsrg} lists first-place TDK papers advised in our group since 2018; \cref{tab:papers-neighbours} lists those from neighbouring faculty groups. Read two or three before writing your spine. Study how they organise the leading figure, contributions and evaluation, regardless of topic. Each link opens the portal page containing the PDF.

{\small
\begin{xltabular}{\linewidth}{>{\raggedright\arraybackslash}p{0.07\linewidth} X >{\raggedright\arraybackslash}p{0.2\linewidth} >{\raggedright\arraybackslash}p{0.24\linewidth}}
  \caption{First-place TDK papers advised in the research group, 2018--2025, from the BME TDK portal.}\label{tab:papers-ftsrg}\\
  \toprule \textbf{Year} & \textbf{Title} & \textbf{Authors} & \textbf{Advisors} \\ \midrule \endfirsthead
  \toprule \textbf{Year} & \textbf{Title} & \textbf{Authors} & \textbf{Advisors} \\ \midrule \endhead
  \bottomrule \endlastfoot
  2018 & \href{https://tdk.bme.hu/conference/VIK/2018/sessions/SW5/paper/Grafmintailleszto-rendszerek-tesztelese}{Functional Testing of Graph Query Engines by Automated Graph Generators} & Bekő Mária & Semeráth Oszkár \\
  2019 & \href{https://tdk.bme.hu/conference/VIK/2019/sessions/intel/paper/Kiberfizikai-rendszerek-megbizhatosagi}{Simulation, inference and verification for state-based advanced reliability models of complex cyber-physical systems} & Nagy Simon József & Marussy Kristóf \\
  2019 & \href{https://tdk.bme.hu/conference/VIK/2019/sessions/model/paper/Tenzorreprezentacios-modszerek-a-komplex}{Towards tensor-based reliability analysis of complex safety-critical systems} & Szekeres Dániel & Marussy Kristóf \\
  2020 & \href{https://tdk.bme.hu/conference/VIK/2020/sessions/modell/paper/Komplex-elosztott-rendszerek}{Towards tensor-based extrafunctional analysis of complex distributed systems} & Szekeres Dániel & Marussy Kristóf \\
  2021 & \href{https://tdk.bme.hu/conference/VIK/2021/sessions/inform/paper/ZKPalapu-audit-digitalis-jegybankpenzt-kezelo}{ZKP-based Audit for Blockchain Systems Managing Central Bank Digital Currency} & Péter Bertalan Zoltán & Kocsis Imre \\
  2022 & \href{https://tdk.bme.hu/conference/VIK/2022/sessions/inform1/paper/Alkalmazasi-szintu-ateresztokepesseg}{Data model driven goodput optimization for execute-order-validate blockchains} & Debreczeni Máté & Klenik Attila, Kocsis Imre \\
  2022 & \href{https://tdk.bme.hu/conference/VIK/2022/sessions/sw8/paper/Kollaborativ-munkafolyamatok-titkossagmegorzo}{Blockchain-based, confidentiality preserving orchestration of collaborative workflows} & Toldi Balázs Ádám & Kocsis Imre \\
  2022 & \href{https://tdk.bme.hu/conference/VIK/2022/sessions/sw8/paper/Graflekerdezesekkel-tamogatott-architektura}{System Architecture Evaluation and Synthesis for Performability with Graph Queries} & Földiák Máté & Marussy Kristóf \\
  2023 & \href{https://tdk.bme.hu/conference/VIK/2023/sessions/inform2/paper/Onkiertekelo-eljarasrendek-tamogatasa}{Self-evaluated policies using Zero-Knowledge proofs} & Farkas Martin, Toldi Balázs Ádám & Kocsis Imre, Péter Bertalan Zoltán \\
  2023 & \href{https://tdk.bme.hu/conference/VIK/2023/sessions/sw9/paper/Valos-ideju-szoftverintenziv-rendszer-modellek}{Abstraction-based model checking for real-time software-intensive system models} & Cziborová Dóra & Marussy Kristóf, Dobos-Kovács Mihály \\
  2024 & \href{https://tdk.bme.hu/conference/VIK/2024/sessions/intel/paper/Kvalitativ-kovetkeztes-alapu-hibadiagnosztika}{Qualitative Reasoning-Based Fault Diagnosis in Distributed Systems} & Tarnay Márton & Földvári András, Péter Bertalan Zoltán \\
  2024 & \href{https://tdk.bme.hu/conference/VIK/2024/sessions/model/paper/Beagyazott-MI-komponensek-rendszer-szintu}{System-level impact analysis of embedded AI components} & Deé-Lukács András Gergely & Pataricza András, Földvári András \\
  2025 & \href{https://tdk.bme.hu/conference/VIK/2025/sessions/inf/paper/Hatasanalizissel-Tamogatott-MI-Biztonsagi}{Impact Assessment Supported AI Safety Wrapper} & Deé-Lukács András Gergely & Pataricza András, Földvári András \\
  2025 & \href{https://tdk.bme.hu/conference/VIK/2025/sessions/inteli/paper/Grafalapu-logikai-megoldok-kombinalasa}{Combining graph-based logic solvers with genetic algorithms} & Tömöri Péter András & Ficsor Attila, Semeráth Oszkár \\
  2025 & \href{https://tdk.bme.hu/conference/VIK/2025/sessions/rendmod/paper/Modellalapu-rendszertervezest-tamogato-agensek}{Developing AI assistants for model-based systems engineering} & Bilicki Vilmos & Elekes Márton, Marussy Kristóf, Vörös András \\
  2025 & \href{https://tdk.bme.hu/conference/VIK/2025/sessions/sw/paper/Numerikus-kovetkezteto-modszerek-vizsgalata}{Numerical constraint propagation for automated graph generation} & Balog Bors & Marussy Kristóf \\
\end{xltabular}

\begin{xltabular}{\linewidth}{>{\raggedright\arraybackslash}p{0.07\linewidth} X >{\raggedright\arraybackslash}p{0.2\linewidth} >{\raggedright\arraybackslash}p{0.24\linewidth}}
  \caption{First-place TDK papers advised in neighbouring groups at the faculty, 2017--2025: the CrySyS Lab (Buttyán Levente) and Pejo Balázs.}\label{tab:papers-neighbours}\\
  \toprule \textbf{Year} & \textbf{Title} & \textbf{Authors} & \textbf{Advisors} \\ \midrule \endfirsthead
  \toprule \textbf{Year} & \textbf{Title} & \textbf{Authors} & \textbf{Advisors} \\ \midrule \endhead
  \bottomrule \endlastfoot
  2017 & \href{https://tdk.bme.hu/conference/VIK/2017/sessions/SW4/paper/a84}{Detection of malicious web pages with static analysis} & Tamás Kristóf & Buttyán Levente \\
  2020 & \href{https://tdk.bme.hu/conference/VIK/2020/sessions/info2/paper/Keresesi-hatekonysag-novelese-nagyon-nagy}{Improving search efficiency in very large malware repositories} & Tamás Csongor & Buttyán Levente \\
  2021 & \href{https://tdk.bme.hu/conference/VIK/2021/sessions/beagy/paper/Kartekony-kod-mentes-allapot-tavoli}{Protocol for remote attestation of malware-free state of embedded IoT devices} & Bak Márton László & Buttyán Levente, Futóné Papp Dorottya \\
  2022 & \href{https://tdk.bme.hu/conference/VIK/2022/sessions/beagy1/paper/Robusztussag-vizsgalat-a-hasonlosagalapu-IoT}{Robustness Against Evasion of Similarity-based IoT Malware Detection Methods} & Sándor József & Buttyán Levente \\
  2023 & \href{https://tdk.bme.hu/conference/VIK/2023/sessions/sw9/paper/Ket-tamadas-a-TLSH-hasonlosagi-hash-algoritmus}{Two Attacks on the TLSH Similarity Digest Scheme} & Fuchs Gábor & Buttyán Levente, Nagy Roland \\
  2024 & \href{https://tdk.bme.hu/conference/VIK/2024/sessions/kvantum/paper/A-securitybydesign-modszer-tamogatasa-az}{Supporting security-by-design via automating attack tree generation from abstract system models} & Szilvágyi Dániel Zoltán & Buttyán Levente \\
  2025 & \href{https://tdk.bme.hu/conference/VIK/2025/sessions/modeling/paper/Nagy-meretu-dinamikus-grafok-klaszterezeset}{Design and analysis of clustering algorithms for large dynamic graphs} & Maliga Dávid & Buttyán Levente \\
  2025 & \href{https://tdk.bme.hu/conference/VIK/2025/sessions/neural/paper/Hozzajarulasi-ertek-szamitasa-elosztott}{Robustness Scores in Federated Learning} & Ulrich Flóra Zsuzsanna & Pejo Balázs \\
\end{xltabular}
}

\section{A thesis from this paper}\label{app:tdk:thesis}
A TDK paper can become a BSc or MSc thesis. The required files are kept as commented inputs in \texttt{thesis.tex}: the thesis title page (\texttt{include/titlepage.tex}), student declaration (\texttt{include/declaration.tex}), faculty structure guide (\texttt{include/guideline.tex}) and task description (\texttt{include/project.tex}). Set \texttt{\textbackslash vikdoktipus} to \texttt{\textbackslash bsc} or \texttt{\textbackslash msc}, then make the following changes from the faculty's guide:
\begin{itemize}[noitemsep]
  \item \textbf{Length.} A BSc thesis is around 30--45 pages, an MSc thesis 60--90 pages of useful content; ask us, departments differ.
  \item \textbf{Language.} The faculty's default is Hungarian. English suits students aiming at international research or industry, but takes more writing effort and more reading effort from advisors and reviewers. This template is English-only. For a Hungarian thesis, use the faculty's template.
  \item \textbf{Format.} A4, 2.5\,cm margins plus 1\,cm for binding, 11--12\,pt serif, one-and-a-half line spacing. The TDK's 15\,mm margins and single spacing are set in \texttt{thesis.tex}; change them there. Typography beyond that: \citet{guide:Jeney}.
  \item \textbf{Structure.} Follow \texttt{include/guideline.tex}: task description, title page, table of contents, declaration, abstracts in both languages, introduction, task analysis, prior work, design and alternatives, evaluation, acknowledgements, bibliography and appendices. A thesis without the mandatory declaration will not be reviewed. The PDF must be at most 15\,MB. The printed copy must include your signed declaration and the topic description signed by the department head, downloaded from the \href{https://diplomaterv.vik.bme.hu/}{dipterv portal}.
\end{itemize}
